% ============================================================================
% Fichier  : partie2/P02_C08_preparer_planifier.tex
% Titre    : Préparer et Planifier la Mission
% Auteur   : MINKA MI NGUIDJOÏ Thierry Emmanuel
% Version  : 1.0 -- novembre 2026
% ============================================================================
% Base : ACSI C06 (quatre questions des termes de référence ; strates
% documentaires ; 21 documents 12/3/6, µ 0,57 U 0,29 H 1,38 ; zones
% retenues stocks, domaine D, Garoua ; continuité reportée).
% Sources : cours Audit_SI de l'auteur (travaux préparatoires, strates,
% lecture du SI en couches, entretiens par niveau), mis en cohérence avec A0
% et A1 ; ISO 19011:2018 (programme d'audit) ; IIA 2024 (planification
% fondée sur les risques).
% Corrigés : corriges/Corr_C08.tex.
% ============================================================================

\chapter{Préparer et Planifier la Mission}
\label{chap:planifier}
\chaptermark{Préparer et planifier}

\epigraph{%
    « Le temps passé à préparer\\
    se reprend sur le terrain, au double. »%
}{Adage des chefs de mission}

\niveauChapitre{Fondamental}{%
    Maîtrise des Chapitres~\ref{chap:risque} (registre, risque d'audit) et
    \ref{chap:mission} (cycle, outils O0 à O9)%
}

\begin{boxobjectifs}
\begin{itemize}
    \item Distinguer les travaux préparatoires de la planification.
    \item Rédiger les termes de référence d'une mission en répondant à quatre
          questions~: quand, comment, par qui, avec quoi.
    \item Établir la liste des documents sollicités, et coter \NEK{} ou
          \NON{} un document absent selon le cas.
    \item Prendre connaissance d'un système d'information par couches, par
          strates documentaires et par niveaux d'interlocuteurs~; conduire
          un entretien de prise de connaissance.
    \item Évaluer a priori le contrôle interne par un questionnaire et un
          test de cheminement.
    \item Choisir les zones d'audit en croisant le risque et la
          connaissance, et rédiger un plan d'audit complet.
\end{itemize}
\end{boxobjectifs}

\bigskip

% ============================================================================
\section*{Question directrice : que faut-il savoir avant d'auditer~?}
% ============================================================================

AUDIX CONSEIL dispose de huit semaines pour auditer LiZbiZ~: quatre
centres de données, douze agences, une application intégrée, des milliers
de comptes et de mouvements de stock. Impossible de tout voir. M.~BEYALA
consacre donc les deux premières semaines à décider \emph{quoi} auditer,
\emph{où}, \emph{avec qui} et \emph{comment}, avant le moindre test.

\medskip
\textbf{Plan du chapitre.}
\begin{enumerate}
    \item Préparer, puis planifier (section~\ref{sec:c08-deux-phases}).
    \item Les termes de référence (section~\ref{sec:c08-tdr}).
    \item Les documents sollicités (section~\ref{sec:c08-documents}).
    \item La prise de connaissance (section~\ref{sec:c08-connaissance}).
    \item L'évaluation a priori et le choix des zones
          (section~\ref{sec:c08-risques}).
    \item Le plan d'audit (section~\ref{sec:c08-plan}).
\end{enumerate}

\separateur

% ============================================================================
\section{Préparer, puis Planifier}
\label{sec:c08-deux-phases}
% ============================================================================

\begin{boxdef}
\termecle{Travaux préparatoires}\index{travaux préparatoires} (phase~0)~:
diligences qui conduisent à la lettre de mission et à la constitution de
l'équipe~: analyse préalable des risques, termes de référence, ressources.

\medskip
\termecle{Planification}\index{planification} (phase~1)~: diligences qui
préparent le terrain à partir de la lettre de mission~: prise de
connaissance, cartographie préliminaire des risques, décomposition de
l'espace d'audit, plan d'audit.
\end{boxdef}

On confond souvent les deux, parce que, du point de vue du contrat, la
mission commence à la signature de la lettre de mission. Du point de vue de
la démarche, la préparation \emph{produit} la lettre de mission~; la
planification l'\emph{exploite}. Les acteurs diffèrent aussi~: la
préparation associe étroitement le commanditaire, la planification
mobilise surtout l'équipe d'audit.

\subsection{L'audit interne~: du plan annuel à la mission}

Pour une fonction d'audit interne, la mission naît d'un plan annuel fondé
sur les risques, approuvé par le comité d'audit~\cite{iia2024gias}. Le plan
couvre l'« univers d'audit » (la liste de tous les objets auditables~:
processus, applications, sites) sur un cycle de trois à cinq ans, avec une
fréquence qui dépend du risque. Pour un cabinet externe comme AUDIX
CONSEIL, la mission naît d'une demande du client~; la préparation commence
par une analyse d'acceptation~: compétence, indépendance, faisabilité dans
le délai.

% ============================================================================
\section{Les Termes de Référence}
\label{sec:c08-tdr}
% ============================================================================

\subsection{Quatre questions}

Les termes de référence (outil~O2) répondent à quatre questions, qui
déterminent le coût et la faisabilité de la mission.

\begin{tableaucours}{L{2.4cm}L{5.0cm}L{4.8cm}}{%
    \tableauHeader{Question} &
    \tableauHeader{Contenu} &
    \tableauHeader{Chez LiZbiZ}
}
Quand~? & Période auditée, durée de la mission & Janvier à avril~2026~;
huit semaines \\
Comment~? & Type d'audit, objectifs, périmètre, référentiels, points
d'attention du commanditaire & Audit d'application (stocks) et audit de
sécurité (contrôle d'accès)~; COSO, ISO/IEC~27002 \\
Par qui~? & Profils requis, auditeurs internes ou externes, absence de
conflit d'intérêts & Un chef de mission, un auditeur réseau et sécurité, un
auditeur données (SQL) \\
Avec quoi~? & Moyens matériels, accès, budget & Accès en lecture à la base,
outil d'analyse de données, déplacements à Garoua et Bafoussam \\
\end{tableaucours}
\vspace{-0.8em}
\begin{center}{\small\itshape Tableau~8.1. Les quatre questions des termes
de référence}\end{center}

\subsection{Écrire un objectif d'audit}

Un objectif d'audit se vérifie~: il dit ce que la mission doit démontrer,
sur quel objet, selon quel critère. « Auditer le système d'information »
n'est pas un objectif~; c'est l'erreur de la première mission.

\begin{tableaucours}{L{5.6cm}L{6.6cm}}{%
    \tableauHeader{Objectif mal écrit} &
    \tableauHeader{Objectif vérifiable}
}
Auditer la sécurité & Évaluer si les accès à l'application intégrée sont
attribués, revus et retirés conformément à la politique de sécurité et à
l'ISO/IEC~27002 (mesures 5.15 à 5.18) \\
Vérifier les stocks & Évaluer si les mouvements de stock de janvier à
avril~2026 sont réels, exhaustifs et correctement valorisés \\
Contrôler Garoua & Évaluer si la facturation de l'agence de Garoua passe
intégralement par l'application intégrée \\
\end{tableaucours}
\vspace{-0.8em}
\begin{center}{\small\itshape Tableau~8.2. Écrire un objectif
d'audit}\end{center}

\subsection{Les livrables de la préparation}

À l'issue de la phase~0, le dossier de la mission contient les termes de
référence, la fiche technique (O1), la lettre de mission signée, et selon
le cas le contrat, l'engagement de confidentialité, la déclaration d'absence
de conflit d'intérêts de chaque auditeur et les autorisations d'accès aux
sites.

% ============================================================================
\section{Les Documents Sollicités}
\label{sec:c08-documents}
% ============================================================================

\subsection{Trois strates de documents}

La liste des documents sollicités (outil~O3) se construit selon le type
d'audit. Les documents n'ont pas tous la même valeur~: on les classe en trois
strates.

\begin{tableaucours}{L{2.4cm}L{3.8cm}L{3.0cm}L{3.0cm}}{%
    \tableauHeader{Strate} &
    \tableauHeader{Exemples} &
    \tableauHeader{Ce qu'ils révèlent} &
    \tableauHeader{Si absents}
}
1. Stratégique & Schéma directeur, politique de sécurité, plan de
continuité, rapports d'audit précédents & L'engagement de la direction &
Gouvernance informelle, recommandations non suivies \\
2. Organisationnelle & Organigramme de la DSI, cartographie des
applications et des interfaces, matrice des droits & Qui fait quoi &
Responsabilités floues, cumuls de fonctions \\
3. Opérationnelle & Journaux d'incidents, de sauvegardes, de traitements~;
inventaires & Ce qui se fait réellement & Impossible de distinguer le
déclaré du réel \\
\end{tableaucours}
\vspace{-0.8em}
\begin{center}{\small\itshape Tableau~8.3. Les strates documentaires}\end{center}

La confrontation des strates est en soi une source de constats~: une
politique de sécurité (strate~1) qui exige une revue trimestrielle des
droits, sans aucune trace de revue dans les journaux (strate~3), révèle
l'écart entre l'affiché et le pratiqué.

\subsection{Un document absent~: \NEK{} ou \NON{}~?}

Chaque document est sollicité avec une priorité (haute ou moyenne), un
interlocuteur et une date de remise. Lorsqu'un document n'est pas remis,
l'auditeur doit qualifier cette absence avec soin.

\begin{tableaucours}{L{4.8cm}cL{5.6cm}}{%
    \tableauHeader{Situation} &
    \tableauHeader{Cote} &
    \tableauHeader{Justification}
}
Le document n'est pas remis malgré une relance, sans qu'on sache s'il
existe & \NEK{} & L'information manque~: on ne peut pas conclure \\
L'audité confirme par écrit que le document n'existe pas, alors que le
critère l'exige & \NON{} & L'absence est établie~: le critère n'est pas
satisfait \\
Le document est remis, mais périmé ou non approuvé & \NON{} & Le critère
(document à jour et approuvé) n'est pas satisfait \\
Le document n'a pas d'objet dans le périmètre (aucune activité concernée)
& sans objet & On le retire du dénominateur, en le justifiant \\
\end{tableaucours}
\vspace{-0.8em}
\begin{center}{\small\itshape Tableau~8.4. Coter un document
manquant}\end{center}

\begin{boxlecture}{Documents sollicités}
Sur 21 documents sollicités le 5~mai, au 22~mai et après une relance~: 12
sont remis et conformes (\OUI{}), 3 sont remis mais périmés (\NON{}), 6 ne
sont pas remis (\NEK{}), dont le plan de continuité et la politique de
sécurité.

\medskip
\textbf{Constatation.} $\mu \approx 0{,}57$~; $U \approx 0{,}29$~;
$H \approx 1{,}38$~bit.

\medskip
\textbf{Analyse} (confiance moyenne). Avant tout test, plus d'un quart de
l'information de base manque, et deux des documents manquants relèvent de
la strate stratégique. Le plan d'audit doit réserver du temps à les
obtenir, ou à établir qu'ils n'existent pas~: chaque \NEK{} levé avant le
terrain réduit l'incertitude de toute la mission.

\medskip
\hfill\lecture{0{,}57}{0{,}29}{1{,}38}
\end{boxlecture}

\begin{boxalert}
\textbf{Violation typique.} Commencer les tests sans avoir obtenu la matrice
des droits, puis conclure sur la séparation des tâches à partir des seuls
entretiens.

\smallskip
\textbf{Règle~:} un document de priorité haute non remis est relancé par
écrit, et son absence est consignée dans le suivi des documents (O9). On ne
conclut jamais sur un contrôle dont la preuve principale manque~: il reste
\NEK{}.
\end{boxalert}

% ============================================================================
\section{Prendre Connaissance du Système d'Information}
\label{sec:c08-connaissance}
% ============================================================================

\subsection{Une prise de connaissance bornée et hiérarchisée}

\begin{boxdef}
\termecle{Prise de connaissance}\index{prise de connaissance}~: ensemble des
diligences par lesquelles l'auditeur acquiert, avant le terrain, une
compréhension suffisante de l'organisation, de son système d'information,
de ses processus et de ses risques pour définir un programme de travail
pertinent. Elle produit le \emph{dossier de prise de connaissance}, première
pièce du dossier d'audit.
\end{boxdef}

Elle est \emph{bornée} par le périmètre de la lettre de mission~: on ne
documente que ce qui le concerne. Elle est \emph{hiérarchisée}~: on lit les
documents par strates. Elle produit un livrable actionnable~: la
cartographie préliminaire des risques. Dans le vocabulaire de
l'auditologie, elle réduit l'incertitude initiale de l'auditeur jusqu'au
niveau qui permet de formuler des hypothèses de risque~\cite{minka2026auditologie}.

\subsection{Lire le SI en couches}

L'auditeur ne peut pas juger une application sans comprendre ce sur quoi
elle repose. Il lit le système d'information en quatre couches, de bas en
haut.

\begin{tableaucours}{L{3.2cm}L{4.6cm}L{4.4cm}}{%
    \tableauHeader{Couche} &
    \tableauHeader{Ce qu'on documente} &
    \tableauHeader{Chez LiZbiZ}
}
1. Infrastructure & Matériels, réseaux, hébergement, liaisons &
Quatre centres de données (Douala, Yaoundé, Garoua, Bafoussam), liaisons
avec les douze agences \\
2. Bases et intergiciels & Systèmes de gestion de bases, interfaces,
échanges & Une base centrale, réplication vers les centres régionaux,
interface avec la banque \\
3. Applications & Modules, paramétrages, contrôles automatisés &
Application intégrée~: achats, stocks, ventes, comptabilité~; tableur
parallèle à Garoua \\
4. Données et accès & Profils, droits, journaux, conservation &
Environ 150~comptes, matrice des droits, journaux d'administration non
protégés \\
\end{tableaucours}
\vspace{-0.8em}
\begin{center}{\small\itshape Tableau~8.5. Le SI de LiZbiZ lu en quatre
couches}\end{center}

\subsection{Interroger par niveau}

Les documents décrivent ce qui est formalisé~; les entretiens révèlent ce
qui est pratiqué. Les entretiens de prise de connaissance ne recueillent
pas encore de preuves~: ils situent les personnes, les tensions et les
problèmes connus.

\begin{tableaucours}{L{3.0cm}L{4.6cm}L{4.6cm}}{%
    \tableauHeader{Niveau} &
    \tableauHeader{Ce qu'on cherche} &
    \tableauHeader{Signal d'alerte}
}
Direction générale & Les objectifs assignés au SI, la tolérance au risque,
l'instance de pilotage & Le dirigeant ignore le rapport d'audit précédent \\
DSI et responsables informatiques & Organisation, parc applicatif,
prestataires, problèmes connus non résolus & « Tout va bien, nous n'avons
aucun incident » \\
Métiers et utilisateurs & La réalité vécue, les contournements, les
incidents non remontés & Un fichier tableur qui double l'application \\
\end{tableaucours}
\vspace{-0.8em}
\begin{center}{\small\itshape Tableau~8.6. Les entretiens de prise de
connaissance}\end{center}

\subsection{Conduire un entretien}

\begin{boxprocedure}
\textbf{Déroulé d'un entretien de prise de connaissance (45 à 60 minutes)}
\begin{enumerate}
    \item \textbf{Ouverture} (5~min)~: rappeler l'objet de la mission, la
          confidentialité, et que l'entretien n'est pas une évaluation de la
          personne.
    \item \textbf{Description} (20~min)~: faire décrire le processus tel
          qu'il est pratiqué, par des questions ouvertes («~Que se passe-t-il
          quand un fournisseur livre~?~»).
    \item \textbf{Approfondissement} (15~min)~: demander les exceptions,
          les incidents, les contournements («~Que faites-vous quand
          l'application est indisponible~?~»).
    \item \textbf{Documents} (5~min)~: noter les documents cités et les
          demander.
    \item \textbf{Problèmes connus} (5~min)~: «~S'il y avait une chose à
          améliorer, laquelle~?~»
    \item \textbf{Compte rendu}~: rédigé dans les 24~heures, daté, et
          transmis à la personne pour correction des faits.
\end{enumerate}
\end{boxprocedure}

\begin{boxalert}
\textbf{Violation typique.} Poser des questions fermées qui suggèrent la
réponse («~Vous revoyez bien les droits chaque trimestre~?~»).

\smallskip
\textbf{Règle~:} en prise de connaissance, on fait décrire~; on ne fait pas
confirmer. Une réponse à une question suggestive ne vaut même pas
déclaration.
\end{boxalert}

\begin{boxmetier}{Auditeur SI}
\textbf{Conduire la prise de connaissance de LiZbiZ}
\begin{enumerate}
    \item Lire d'abord les documents de strate~1, puis~2, puis~3, et noter
          chaque écart entre strates.
    \item Rencontrer le directeur général, puis le responsable informatique,
          puis deux chefs d'agence et un magasinier.
    \item Demander à chacun les « problèmes connus non résolus »~: ce sont
          souvent les risques les plus sérieux.
    \item Cartographier les couches du SI, du centre de données au compte
          utilisateur.
    \item Consigner chaque entretien dans un compte rendu daté.
\end{enumerate}
\end{boxmetier}

% ============================================================================
\section{Évaluer a Priori et Choisir les Zones}
\label{sec:c08-risques}
% ============================================================================

\subsection{L'évaluation a priori du contrôle interne}

Avant les tests, l'auditeur forme une première opinion sur le contrôle
interne, à partir des documents et des entretiens~: les contrôles
existent-ils, sont-ils conçus pour traiter les risques
(Chapitre~\ref{chap:controle-interne})~? Cette évaluation ne conclut rien~;
elle oriente. Un contrôle jugé bien conçu sera testé en efficacité~; un
contrôle absent ou mal conçu n'a pas besoin d'être testé~: c'est déjà un
constat, et l'auditeur portera ses tests directement sur les données.

Deux outils la structurent~:
\begin{itemize}
    \item le \termecle{questionnaire de contrôle interne}\index{questionnaire
          de contrôle interne}, qui liste les contrôles attendus pour chaque
          processus et demande s'ils existent~; ses réponses sont des
          déclarations, à vérifier~;
    \item le \textbf{test de cheminement} (Chapitre~\ref{chap:controle-interne}),
          qui suit une opération réelle et vérifie que les contrôles déclarés
          existent bien.
\end{itemize}

\begin{boxscenario}{Opération LIZBIZ. Un cheminement de réception}
\textbf{Faits.} L'auditeur suit la commande CA-2026-0412 (200~ordinateurs
portables, Douala)~: bon de commande validé par le DAF le 6~février~;
réception saisie par M.~FOUDA le 19~février~; facture rapprochée par
Mme~OWONA le 24~février~; paiement validé par Mme~MBALLA le 2~mars.

\medskip
\textbf{Constatation.} Le rapprochement automatique commande, réception,
facture existe et a fonctionné sur cette opération. En revanche, M.~FOUDA a
saisi le même jour un ajustement négatif de 12~unités sur la même référence,
sans pièce.

\medskip
\textbf{Analyse.} Le contrôle de rapprochement est bien conçu et sera testé
en efficacité~; le contrôle des ajustements, absent, est déjà un constat
(Chapitre~\ref{chap:fraude}).
\end{boxscenario}

\subsection{Le risque et la connaissance}

L'auditeur part de la cartographie des risques de l'organisation, si elle
existe (Chapitre~\ref{chap:risque}). Sinon, il en établit une à partir de
la prise de connaissance. Il y ajoute ce que l'auditologie invite à
mesurer~: ce que l'on sait de chaque risque. On audite en priorité ce qui
est \emph{le plus risqué} et \emph{le moins bien connu}.

\noindent\begin{minipage}{\linewidth}
\begin{center}
\begin{tikzpicture}[scale=1.0]
\fill[orange!25] (0,2) rectangle (3,4);
\fill[red!30] (3,2) rectangle (6,4);
\fill[green!20] (0,0) rectangle (3,2);
\fill[orange!15] (3,0) rectangle (6,2);
\draw[thick] (0,0) rectangle (6,4);
\draw (3,0) -- (3,4); \draw (0,2) -- (6,2);
\node[font=\scriptsize, align=center] at (1.5,3) {Risque élevé, connu\\ tester en priorité};
\node[font=\scriptsize\bfseries, align=center] at (4.5,3) {Risque élevé, peu connu\\ zone prioritaire};
\node[font=\scriptsize, align=center] at (1.5,1) {Risque faible, connu\\ revue légère};
\node[font=\scriptsize, align=center] at (4.5,1) {Risque faible, peu connu\\ réduire l'incertitude};
\node[font=\small] at (3,-0.45) {Incertitude sur le risque (faible à forte)};
\node[font=\small, rotate=90] at (-0.45,2) {Risque résiduel};
\node[font=\scriptsize\bfseries, color=CouleurPrimaire] at (5.2,3.55) {D};
\node[font=\scriptsize\bfseries, color=CouleurPrimaire] at (2.2,3.55) {Stocks};
\node[font=\scriptsize\bfseries, color=CouleurPrimaire] at (4.0,3.55) {Garoua};
\node[font=\scriptsize\bfseries, color=CouleurPrimaire] at (5.0,1.5) {Continuité};
\end{tikzpicture}
\end{center}
\vspace{-0.4em}
\begin{center}{\small\itshape Figure~8.1. Choisir les zones d'audit~: le
risque et la connaissance}\end{center}
\end{minipage}

\begin{boxscenario}{Opération LIZBIZ. Choisir les zones d'audit}
\textbf{Constatations.} Le registre de Mme~NKENG classe trois risques comme
intolérables (ajustement du stock, fournisseurs fictifs, comptes d'agents
partis). La liste des documents laisse inconnus le plan de continuité et la
politique de sécurité. Les entretiens signalent un fichier tableur de suivi
des stocks tenu à Garoua en parallèle de l'application.

\medskip
\textbf{Analyse} (confiance moyenne). Le plan d'audit retient trois zones~:
la gestion des stocks (risque élevé, connaissance moyenne), le contrôle
d'accès, domaine~D (risque élevé, connaissance faible), et l'agence de
Garoua (contournement de l'application). La continuité est reportée à une
mission ultérieure, faute de temps, et le rapport le dira.
\end{boxscenario}

\begin{boxinfo}
\textbf{Reporter n'est pas ignorer.} La continuité est placée dans la
région « risque faible, peu connu » sur la seule foi de la cotation de R4
(Chapitre~\ref{chap:risque}). Le plan de continuité n'ayant pas été remis,
cette position est elle-même incertaine. Le plan d'audit l'écrit~: «~zone
non auditée, plan de continuité non remis (\NEK{}), à auditer en
priorité lors de la mission suivante~».
\end{boxinfo}

% ============================================================================
\section{Rédiger le Plan d'Audit}
\label{sec:c08-plan}
% ============================================================================

\subsection{Dimensionner les tests}

L'étendue des tests dépend de la fréquence du contrôle testé et du risque.
Les ordres de grandeur du Tableau~3.9 (Chapitre~\ref{chap:controle-interne})
s'appliquent~; on retient le haut de la fourchette lorsque le risque est
élevé ou le contrôle critique, le bas lorsqu'il est faible. Lorsque le
risque de détection doit être très bas, comme pour les stocks de LiZbiZ
(RD~$\approx$~0,07, Chapitre~\ref{chap:risque}), on complète l'échantillon
par des tests sur l'ensemble des données. Le Chapitre~\ref{chap:preuve}
traite de l'échantillonnage statistique.

\subsection{Le contenu du plan}

Le plan d'audit (outil~O6), ou rapport d'orientation, clôt la
planification. Il est la feuille de route du terrain.

\begin{boxprocedure}
\textbf{Le contenu du plan d'audit.} Pour qu'il puisse être approuvé, puis
servir de référence lors de la revue de la mission, le plan d'audit
indique~:
\begin{itemize}
    \item le contexte, les objectifs et le périmètre, rappelés de la lettre
          de mission~;
    \item les référentiels et les critères retenus~;
    \item l'évaluation a priori du contrôle interne et la cartographie
          préliminaire des risques~;
    \item les zones d'audit retenues et celles écartées, avec la raison~;
    \item le programme de travail par zone~: domaines décomposés, tests,
          tailles d'échantillon, preuves attendues~;
    \item le calendrier, l'équipe et l'affectation des tâches~;
    \item la date et la signature de son approbation par le commanditaire.
\end{itemize}
\end{boxprocedure}

\subsection{Le programme de travail de LiZbiZ}

\begin{tableaucours}{L{2.4cm}L{4.4cm}L{2.8cm}L{2.6cm}}{%
    \tableauHeader{Zone} &
    \tableauHeader{Tests principaux} &
    \tableauHeader{Étendue} &
    \tableauHeader{Auditeur}
}
Stocks & Ajustements par magasinier et par jour~; rapprochement
réceptions et commandes~; inventaire contradictoire de 50~références &
Population complète (requêtes)~; inventaire sur 4~sites & Auditeur
données \\
Domaine~D & Cycle de vie des comptes, revue des droits, comptes à
privilèges, journalisation & 31~départs (population)~; 25~comptes à
privilèges & Auditeur sécurité \\
Garoua & Rapprochement tableur et application~; entretien du chef
d'agence & Toute l'année~2025 & Chef de mission \\
Paiements (ajouté) & Paiements sous le seuil, RIB partagés & Population
complète & Auditeur données \\
\end{tableaucours}
\vspace{-0.8em}
\begin{center}{\small\itshape Tableau~8.7. Programme de travail de la
mission LiZbiZ (extrait)}\end{center}

\begin{boiteRemarque}{Le plan d'audit peut-il changer en cours de mission~?}
Oui, et il le doit lorsque le terrain révèle un risque imprévu~: les
ajustements de fin de mois à Douala et les paiements de TRANSLOG SERVICES
(Chapitre~\ref{chap:fraude}) n'étaient pas dans le plan initial d'AUDIX
CONSEIL. Mais toute modification est écrite, motivée et approuvée par le
commanditaire~; sinon, le périmètre réellement audité ne correspond plus à
celui que le rapport annonce.
\end{boiteRemarque}

\begin{boxmetier}{Contrôleur interne}
\textbf{Préparer l'arrivée des auditeurs.} Mme~NKENG centralise les
demandes de la liste O3, fixe à chaque direction une date de remise et tient
le tableau de suivi. Elle signale au chef de mission, dès la première
semaine, les documents qui n'existent pas~: un \NON{} établi tôt vaut mieux
qu'un \NEK{} découvert en fin de terrain.
\end{boxmetier}

\begin{boxmetier}{Gestionnaire de risques}
\textbf{Mettre le registre à la disposition de l'audit.} Le registre des
risques, avec ses justifications, est l'une des premières pièces remises.
L'auditeur ne refait pas la cartographie~; il l'utilise, et la teste.
\end{boxmetier}

\begin{boxinfo}
\textbf{Ce que dit la recherche.} L'article A0~\cite{minka2026auditologie}
décrit la mission comme une suite de décisions sous incertitude. La
planification en est la première~: elle choisit où dépenser l'effort de
preuve. Le critère « le plus risqué et le moins connu » revient à placer
l'effort là où il réduit le plus l'incertitude qui compte, c'est-à-dire
celle qui porte sur des risques élevés.
\end{boxinfo}

\separateur

\begin{boxresume}
\textbf{Ce qu'il faut retenir de ce chapitre~:}
\begin{itemize}
    \item \textbf{La préparation produit la lettre de mission~; la
          planification l'exploite.}
    \item \textbf{Les termes de référence} répondent à quatre questions~:
          quand, comment, par qui, avec quoi~; un objectif d'audit se
          vérifie.
    \item \textbf{Les documents se lisent par strates}~: stratégique,
          organisationnelle, opérationnelle~; leur confrontation révèle
          l'écart entre l'affiché et le pratiqué.
    \item \textbf{Un document absent} est coté \NEK{} si l'on ne sait pas
          s'il existe, \NON{} si son absence est établie alors que le
          critère l'exige.
    \item \textbf{La prise de connaissance} lit le SI en quatre couches,
          interroge par niveaux, et fait décrire plutôt que confirmer.
    \item \textbf{L'évaluation a priori} (questionnaire, cheminement)
          oriente~; un contrôle absent est déjà un constat.
    \item \textbf{Le plan d'audit} retient les zones les plus risquées et les
          moins connues, dit ce qui est écarté, dimensionne les tests et se
          fait approuver~; il évolue par écrit.
\end{itemize}
\end{boxresume}

\begin{boxinfo}
\textbf{Points de contrôle de l'élève.} Je sais~: écrire un objectif
d'audit vérifiable~; classer un document par strate~; coter un document
manquant~; décrire un SI en quatre couches~; conduire un entretien en six
temps~; placer une zone dans la grille risque et connaissance.
\end{boxinfo}

\bigskip

\begin{boxexo}[title={\ding{45}\quad Exercice~8.1. Classer les documents \niveaubase}]
Classez dans l'une des trois strates les documents suivants, et indiquez ce
que révélerait leur absence~: le schéma directeur, la matrice des droits, le
journal des sauvegardes, le rapport de la dernière mission d'audit,
l'inventaire des licences, la charte d'utilisation du SI.
\end{boxexo}

\begin{boxexo}[title={\ding{45}\quad Exercice~8.2. Coter les absences \niveaubase}]
Pour chacune des situations suivantes, indiquez la cote du contrôle « un
plan de continuité approuvé existe pour le centre de Garoua », et
justifiez.
\begin{enumerate}[(a)]
    \item Le plan est remis~; il date de 2019 et n'a jamais été approuvé.
    \item Le responsable informatique dit qu'il existe, mais ne le remet pas
          malgré deux relances.
    \item Le directeur général écrit qu'aucun plan n'a jamais été rédigé.
    \item Le plan est remis, approuvé en 2025, et couvre Garoua.
\end{enumerate}
\end{boxexo}

\begin{boxexo}[title={\ding{45}\quad Exercice~8.3. Écrire des objectifs \niveaumoyen}]
Réécrivez en objectifs vérifiables~: (a) «~Auditer la paie de l'ANPP~»~;
(b) «~Voir si la banque est sécurisée~»~; (c) «~Contrôler le prestataire
DIGICORE~».
\end{boxexo}

\begin{boxexo}[title={\ding{45}\quad Exercice~8.4. Le cas des documents \niveaumoyen}]
Deux des six documents non remis par LiZbiZ sont finalement remis le
29~mai, conformes~; un troisième est déclaré inexistant par écrit.
Recalculez $\mu$, $U$ et $H$ (posez le calcul) et commentez.
\end{boxexo}

\begin{boxexo}[title={\ding{45}\quad Exercice~8.5. Une question suggestive \niveaumoyen}]
Réécrivez en questions ouvertes de prise de connaissance~: (a) «~Les
sauvegardes sont faites tous les jours~?~»~; (b) «~Vous avez un plan de
continuité~?~»~; (c) «~Les droits des partants sont bien retirés~?~»
\end{boxexo}

\begin{boxexo}[title={\ding{45}\quad Exercice~8.6. Les couches de la BANQUE SANAGA \niveaumoyen}]
Décrivez le SI de la BANQUE SANAGA en quatre couches, à partir de ce que
vous savez du cas (système central, agences, prestataire DIGICORE). Pour
chaque couche, indiquez un document à solliciter.
\end{boxexo}

\begin{boxexo}[title={\ding{45}\quad Exercice~8.7. Un cheminement \niveaumoyen}]
Décrivez le test de cheminement d'un paiement fournisseur chez LiZbiZ~:
les étapes, les documents examinés à chaque étape, les contrôles attendus.
Qu'en déduisez-vous si, à l'étape du paiement, le valideur est la personne
qui a créé le fournisseur~?
\end{boxexo}

\begin{boxexo}[title={\ding{45}\quad Exercice~8.8. Placer les zones \niveaumoyen}]
Placez dans la grille de la Figure~8.1 les zones suivantes de la BANQUE
SANAGA, et justifiez~: la gestion des changements du système central, la
messagerie, le contrat DIGICORE, la paie interne.
\end{boxexo}

\begin{boxexo}[title={\ding{45}\quad Exercice~8.9. Le plan d'audit de l'ANPP \niveauexpert}]
La Chambre des comptes prépare l'audit des systèmes d'information de l'ANPP
(Chapitre~\ref{chap:audit-auditologie}).
\begin{enumerate}[(a)]
    \item Rédigez les termes de référence en répondant aux quatre questions.
    \item Établissez la liste des dix documents à solliciter en priorité,
          classés par strate.
    \item Proposez trois zones d'audit et justifiez-les par le risque et par
          l'état de la connaissance.
    \item Dimensionnez l'échantillon du test de la revue mensuelle des droits
          d'accès au système de paiement.
\end{enumerate}
\end{boxexo}

\begin{boxexo}[title={\ding{45}\quad Exercice~8.10. Modifier le plan \niveauexpert}]
Le 3~juin, l'auditeur données découvre les paiements de TRANSLOG SERVICES.
Rédigez la note de modification du plan d'audit adressée au commanditaire~:
motif, zone ajoutée, tests, effet sur le calendrier et le budget, ce qui est
retiré pour compenser.
\end{boxexo}

\bigskip

\begin{boxinfo}
\textbf{Pour aller plus loin.}
\begin{itemize}
    \item ISO~19011:2018~\cite{iso19011_2018}~: la préparation des
          activités d'audit et la revue documentaire.
    \item Le guide du CHAI~\cite{chai2014guide}~: des listes de documents
          par type d'audit.
    \item IIA, Normes mondiales~\cite{iia2024gias}~: la planification
          fondée sur les risques.
\end{itemize}
\end{boxinfo}

\begin{boxinfo}
\textbf{Transition.} Le plan d'audit a désigné les zones à auditer. Il faut
maintenant les découper en éléments que l'auditeur peut réellement
vérifier. Le Chapitre~\ref{chap:decomposition} présente le gabarit canonique
de décomposition de l'espace d'audit et la fiche d'assertion.
\end{boxinfo}
